% ==========================================================================
%  NOMENCLATURA  —  colocar en el front matter (antes de la Introducción)
%  Requiere: \usepackage{longtable}
%  Tablas agrupadas por categoría. Adaptar/completar según capítulos finales.
% ==========================================================================

\chapter*{Nomenclatura}
\addcontentsline{toc}{chapter}{Nomenclatura}
\markboth{Nomenclatura}{Nomenclatura}

\renewcommand{\arraystretch}{1.25}

% --------------------------------------------------------------------------
\section*{Números adimensionales}
\begin{longtable}{@{}p{0.16\textwidth} p{0.78\textwidth}@{}}
$Fr$        & Número de Froude, $Fr = V/\sqrt{g\,L_{WL}}$ \\
$Fr_M,\,Fr_S$ & Número de Froude en escala modelo y buque \\
$Fr_{tr}$   & Número de Froude de la popa espejo, $Fr_{tr} = V/\sqrt{g\,T_{tr}}$ \\
$Re$        & Número de Reynolds, $Re = V L_{OS}/\nu$ \\
$Re_M,\,Re_S$ & Número de Reynolds en escala modelo y buque \\
$\overline{Re_M}$ & Número de Reynolds promedio en escala modelo (correlación de $k_{tr}$) \\
$Re_n$      & Número de Reynolds rotacional de la hélice, $Re_n = n D^2/\nu$ \\
$Re_{n,\mathrm{ref}}$ & Número de Reynolds rotacional de la escala de referencia \\
$Re_c$      & Número de Reynolds seccional de la hélice, $Re_c = c_{r/R}\,V_{\text{ref}}(r)/\nu$; se indica la sección empleada como $Re_{c,0.7R}$ o $Re_{c,0.75R}$ \\
$Re_{\theta t}$ & Número de Reynolds de transición basado en el espesor de cantidad de movimiento, obtenido de la correlación empírica en la corriente libre \\
$\widetilde{Re}_{\theta t}$ & Variable de transporte correspondiente a $Re_{\theta t}$ en el modelo de transición SST-LM \\
$y^+$       & Distancia adimensional a la pared, $y^+ = u_\tau y/\nu$ (primer centro de celda) \\
$Co,\,Co_{max}$ & Número de Courant, $Co = |\mathbf{U}|\,\Delta t/\Delta x$, y su valor máximo \\
$Tu$        & Intensidad turbulenta, $Tu = \sqrt{2k/3}\,/\,|\mathbf{U}|$, evaluada en el plano del disco de la hélice \\
$C_B$       & Coeficiente de bloque \\
$C_P$       & Coeficiente de presión sobre el casco, $C_P = (p-p_\infty)/\tfrac{1}{2}\rho V^2$ \\
$C_p$       & Coeficiente de presión sobre la pala, $C_p = (p-p_\infty)/\tfrac{1}{2}\rho V_{\text{ref}}^2$, con $V_{\text{ref}}(r)$ la velocidad resultante en la sección de radio $r$ \\
$\Delta C_p$ & Diferencia del coeficiente de presión sobre la pala entre escalas, $\Delta C_p = C_p(Re_n) - C_p(Re_{n,\mathrm{ref}})$ \\
$C_f$       & Coeficiente de fricción local sobre la pala, $C_f = |\boldsymbol{\tau}_w|/\tfrac{1}{2}\rho V_{\text{ref}}^2$ \\
\end{longtable}

\noindent{\footnotesize Los coeficientes referidos al casco se denotan con subíndice mayúscula ($C_P$, $C_F$) y los referidos a la pala del propulsor con subíndice minúscula ($C_p$, $C_f$).\par}

% --------------------------------------------------------------------------
\section*{Coeficientes y fuerzas de resistencia}
\begin{longtable}{@{}p{0.16\textwidth} p{0.78\textwidth}@{}}
$C_T$              & Coeficiente de resistencia total \\
$C_{TM},\,C_{TS}$  & Coeficiente de resistencia total en escala modelo y buque \\
$C_F$              & Coeficiente de resistencia friccional \\
$C_{F0}$           & Coeficiente de fricción de referencia (línea ITTC-57) \\
$C_{FM0},\,C_{FS0}$& Coeficiente de fricción de referencia en escala modelo y buque \\
$C_V$              & Coeficiente de resistencia viscosa, $C_V = C_F + C_{PV}$ \\
$C_{PV}$           & Coeficiente de resistencia por presión viscosa \\
$C_W$              & Coeficiente de resistencia por formación de olas \\
$C_R$              & Coeficiente de resistencia residual (componente no friccional en la descomposición de Hughes/Froude) \\
$\Delta C_F$       & Corrección por rugosidad de la superficie (Townsin \& Dey) \\
$C_A$              & Coeficiente de correlación modelo--buque \\
$C_{AAS}$          & Coeficiente de resistencia aerodinámica de la obra muerta \\
$C_{tr,S}$         & Coeficiente de resistencia asociado al espejo de popa a escala buque \\
$R_T$              & Fuerza de resistencia total \\
$R_{TM},\,R_{TS}$  & Fuerza de resistencia total en escala modelo y buque \\
$R_V$              & Resistencia viscosa \\
$R_F$              & Resistencia friccional \\
$R_{PV}$           & Resistencia por presión viscosa \\
$R_W$              & Resistencia por formación de olas \\
$k,\;(1+k)$        & Factor de forma \\
$k_M,\,k_S$        & Factor de forma en escala modelo y buque \\
$k_{tr}$           & Factor de forma del espejo \\
$V$                & Velocidad de avance del buque o modelo \\
$V_M,\,V_S$        & Velocidad de avance en escala modelo y buque \\
$q$                & Presión dinámica, $q = \tfrac{1}{2}\rho V^2$ \\
$P_{ES}$           & Potencia efectiva a escala buque, $P_{ES} = R_{TS}V_S$ \\
\end{longtable}

% --------------------------------------------------------------------------
\section*{Geometría del casco y del espejo}
\begin{longtable}{@{}p{0.16\textwidth} p{0.78\textwidth}@{}}
$L_{WL}$    & Eslora en flotación \\
$L_{PP}$    & Eslora entre perpendiculares \\
$L_{OS}$    & Eslora total sumergida \\
$B$         & Manga \\
$T$         & Calado \\
$S$         & Área de superficie mojada \\
$\nabla$    & Volumen de desplazamiento \\
$\Delta$    & Desplazamiento (masa) \\
$\lambda$   & Razón de escala geométrica \\
$L/B$       & Relación de esbeltez (eslora--manga) \\
$k_s$       & Altura de rugosidad equivalente \\
$A_{tr}$    & Área del espejo sumergido \\
$A_{\max}$  & Área máxima de la sección transversal \\
$T_{tr}$    & Inmersión efectiva del borde inferior del espejo \\
$y_{tr}$    & Semimanga en la intersección del espejo con la superficie libre \\
$tr_{\text{ratio}}$    & Relación entre el área del espejo sumergido y la sección maestra, $A_{tr}/A_{\max}$ \\
$tr_{\text{fullness}}$ & Coeficiente de llenado del espejo, $A_{tr}/(2\,y_{tr}\,T_{tr})$ \\
\end{longtable}

% --------------------------------------------------------------------------
\section*{Propulsión (hélices en aguas abiertas)}
\begin{longtable}{@{}p{0.16\textwidth} p{0.78\textwidth}@{}}
$K_T$       & Coeficiente de empuje, $K_T = T/(\rho n^2 D^4)$ \\
$K_{TM},\,K_{TS}$ & Coeficiente de empuje en escala modelo y buque \\
$K_Q$       & Coeficiente de torque, $K_Q = Q/(\rho n^2 D^5)$ \\
$K_{QM},\,K_{QS}$ & Coeficiente de torque en escala modelo y buque \\
$K_{Tp},\,K_{Tf}$ & Componentes de presión y de fricción del coeficiente de empuje \\
$K_{Qp},\,K_{Qf}$ & Componentes de presión y de fricción del coeficiente de torque \\
$\Delta K_T$      & Incremento del coeficiente de empuje entre escalas, $\Delta K_T = K_{TS}-K_{TM}$. El procedimiento ITTC-78 emplea la convención de signo opuesta \\
$\Delta K_Q$      & Corrección de escala sobre el coeficiente de torque \\
$\Delta K_{Tp},\,\Delta K_{Tf}$ & Componentes de presión y de fricción del efecto de escala sobre $K_T$ \\
$\Delta K_{Qp},\,\Delta K_{Qf}$ & Componentes de presión y de fricción del efecto de escala sobre $K_Q$ \\
$\delta K_T,\,\delta K_Q$ & Diferencia entre el cálculo sensible a la transición y el totalmente turbulento, a escala e igual punto de funcionamiento fijos \\
$\delta K_{Tp},\,\delta K_{Tf}$ & Componentes de presión y de fricción de $\delta K_T$ \\
$\delta K_{Qp},\,\delta K_{Qf}$ & Componentes de presión y de fricción de $\delta K_Q$ \\
$f_p$       & Fracción de presión del efecto de escala, $f_p = \Delta K_{Tp}/\Delta K_T$ \\
$J$         & Coeficiente de avance, $J = V_a/(nD)$ \\
$\eta_0$    & Rendimiento en aguas abiertas \\
$\eta_{0M},\,\eta_{0S}$ & Rendimiento en aguas abiertas en escala modelo y buque \\
$T$         & Empuje \\
$T_p,\,T_f$ & Componentes de presión y de fricción del empuje \\
$Q$         & Torque \\
$Q_p,\,Q_f$ & Componentes de presión y de fricción del torque \\
$n$         & Velocidad de giro [rev/s] \\
$D$         & Diámetro de la hélice \\
$R$         & Radio de la pala, $R = D/2$ \\
$r$         & Coordenada radial medida desde el eje \\
$r/R$       & Posición radial adimensional \\
$V_a$       & Velocidad de avance del propulsor, $V_a = J\,n\,D$ \\
$V_{\text{ref}}(r)$ & Velocidad resultante del flujo incidente en la sección de radio $r$, $V_{\text{ref}}(r)=\sqrt{V_a^2+(2\pi n r)^2}$ \\
$V_{0.7R}$  & Velocidad resultante en la sección $r/R=0{,}7$ \\
$Z$         & Número de palas \\
$A_E,\,A_0$ & Área expandida de las palas y área del disco \\
$A_E/A_0$   & Relación de área expandida \\
$A_b$       & Superficie del álabe (dominio de integración de fuerzas) \\
$P$         & Paso geométrico \\
$P/D$       & Relación paso--diámetro \\
$c_{r/R}$   & Cuerda de la sección en la posición radial $r/R$ \\
$c_{0.7R},\,c_{0.75R}$ & Cuerda de la sección a $r/R=0{,}7$ y $r/R=0{,}75$ \\
$t_{\max}$  & Espesor máximo de la sección de pala \\
$k_{s,p}$   & Rugosidad equivalente de la pala (ITTC-78) \\
$x/c$       & Coordenada adimensional a lo largo de la cuerda de la sección \\
$C_{D,M},\,C_{D,S}$ & Coeficiente de arrastre de perfil en escala modelo y buque (ITTC-78) \\
$\Delta C_D$& Corrección de arrastre de perfil (ITTC-78) \\
\end{longtable}

% --------------------------------------------------------------------------
\section*{Modelado numérico y turbulencia}
\begin{longtable}{@{}p{0.16\textwidth} p{0.78\textwidth}@{}}
$U_i,\,\mathbf{U}$ & Componente $i$ de la velocidad media y su forma vectorial \\
$u_i'$      & Fluctuación turbulenta de la velocidad \\
$x_i$       & Coordenada espacial \\
$x,\,y,\,z$ & Coordenadas longitudinal, transversal y vertical; $z$ positiva hacia arriba desde la flotación estática. En capa límite, $y$ denota la distancia normal a la pared \\
$t$         & Tiempo \\
$\Delta t,\,\Delta x$ & Paso temporal y tamaño local de celda \\
$p,\,p_\infty$ & Presión estática local y presión de la corriente no perturbada \\
$p_{rgh}$   & Presión modificada, $p_{rgh}=p-\rho g_j x_j$ \\
$g,\,g_i$   & Aceleración de la gravedad y su componente $i$ \\
$\alpha_{water}$ & Fracción volumétrica de agua (VoF) \\
$\rho$      & Densidad del fluido \\
$\nu$       & Viscosidad cinemática molecular \\
$\nu_t$     & Viscosidad cinemática turbulenta \\
$\mu,\,\mu_t$ & Viscosidad dinámica molecular y turbulenta \\
$k$         & Energía cinética turbulenta \\
$\omega$    & Tasa específica de disipación \\
$\varepsilon$ & Tasa de disipación turbulenta \\
$k_L$       & Energía cinética de las fluctuaciones laminares (modelo $k$-$k_L$-$\omega$) \\
$\tau_w,\,\boldsymbol{\tau}_w$ & Esfuerzo cortante de pared y su forma vectorial \\
$u_\tau$    & Velocidad de fricción, $u_\tau=\sqrt{\tau_w/\rho}$ \\
$\mathbf{n}$ & Versor normal exterior a la superficie \\
$\hat{\mathbf{e}}_x$ & Versor en la dirección de avance \\
$\delta_{ij}$ & Delta de Kronecker \\
$\gamma$    & Intermitencia (variable de transporte del modelo de transición SST-LM) \\
$F_1,\,F_2$ & Funciones de mezcla del modelo SST \\
$\tilde{P}_k$ & Término de producción limitado de turbulencia \\
$S_{ij},\,S$& Tensor y magnitud de la tasa de deformación \\
$a_1$       & Constante de Bradshaw ($=0{,}31$) \\
$\alpha,\,\beta,\,\beta^*$ & Coeficientes de cierre del modelo $k$--$\omega$ SST \\
$\sigma_k,\,\sigma_\omega,\,\sigma_{\omega 2}$ & Coeficientes de difusión del modelo $k$--$\omega$ SST \\
$C_\mu$     & Constante del modelo $k$--$\varepsilon$ ($=0{,}09$) \\
$C_{\varepsilon 1},\,C_{\varepsilon 2}$ & Constantes del modelo $k$--$\varepsilon$ \\
$\sigma_\varepsilon$ & Coeficiente de difusión de $\varepsilon$ \\
$Q$         & Segundo invariante del gradiente de velocidad (criterio Q) \\
\end{longtable}

% --------------------------------------------------------------------------
\section*{Verificación, validación e incertidumbre}
\begin{longtable}{@{}p{0.16\textwidth} p{0.78\textwidth}@{}}
$GCI$       & Índice de convergencia de malla (\textit{Grid Convergence Index}) \\
$R$         & Razón de convergencia \\
$p$         & Orden aparente de precisión \\
$r$         & Razón de refinamiento de malla \\
$h$         & Tamaño característico de celda \\
$\Phi$      & Variable de interés en el estudio de convergencia \\
$\epsilon$  & Diferencia de la solución entre mallas sucesivas \\
$F_s$       & Factor de seguridad ($=1{,}25$) \\
$u_i$       & Componente $i$ de incertidumbre estándar \\
$u_c$       & Incertidumbre estándar combinada \\
$k_c$       & Factor de cobertura ($=2$) \\
$U_{95}$    & Incertidumbre expandida (nivel de confianza del 95\,\%) \\
$N$         & Número de repeticiones de un ensayo o de puntos de comparación \\
$E\%D$      & Error de comparación, $E\%D = (D-S)/D\times100$ \\
$D$         & Valor experimental (\textit{data}) \\
$S$         & Resultado computacional (\textit{simulation}) \\
$U_D$       & Incertidumbre experimental \\
$U_{num}$   & Incertidumbre numérica \\
$U_{val}$   & Incertidumbre de validación, $U_{val}=\sqrt{U_D^2+U_{num}^2}$ \\
MAE         & Error absoluto medio (\textit{Mean Absolute Error}) \\
\end{longtable}

% --------------------------------------------------------------------------
\section*{Acrónimos}
\begin{longtable}{@{}p{0.16\textwidth} p{0.78\textwidth}@{}}
CFD     & \textit{Computational Fluid Dynamics} (dinámica de fluidos computacional) \\
EFD     & \textit{Experimental Fluid Dynamics} (dinámica de fluidos experimental) \\
RANS    & \textit{Reynolds-Averaged Navier--Stokes} \\
SST     & \textit{Shear Stress Transport} (modelo de turbulencia) \\
SST-LM  & Modelo de transición SST de Langtry--Menter ($\gamma$-$\widetilde{Re}_{\theta t}$) \\
EASM    & \textit{Explicit Algebraic Stress Model} (modelo de esfuerzos algebraico explícito) \\
VoF     & \textit{Volume of Fluid} \\
MULES   & \textit{Multidimensional Universal Limiter for Explicit Solution} (compresión de interfaz en VoF) \\
FVM     & Método de volúmenes finitos (\textit{Finite Volume Method}) \\
SIMPLE  & \textit{Semi-Implicit Method for Pressure-Linked Equations} \\
PISO    & \textit{Pressure-Implicit with Splitting of Operators} \\
PIMPLE  & Algoritmo híbrido PISO--SIMPLE \\
MRF     & Marco de referencia múltiple (\textit{Multiple Reference Frame}) \\
AMI     & Interfaz de malla arbitraria (\textit{Arbitrary Mesh Interface}) \\
DB      & Doble cuerpo (\textit{double body}) \\
NFL     & Línea de fricción numérica (\textit{Numerical Friction Line}) \\
ITTC    & \textit{International Towing Tank Conference} \\
ITTC-57 & Línea de correlación modelo--buque de la ITTC (1957) \\
ITTC-78 & Método de predicción de potencia de la ITTC (1978) \\
GCI     & Índice de convergencia de malla \\
V\&V    & Verificación y validación \\
MAE     & Error absoluto medio (\textit{Mean Absolute Error}) \\
RSS     & Raíz de la suma de los cuadrados (\textit{Root Sum Squares}) \\
RMS     & Raíz del valor cuadrático medio (\textit{Root Mean Square}) \\
POW     & Ensayo de hélice en aguas abiertas (\textit{Propeller Open Water}) \\
BPG     & Guías de buenas prácticas (\textit{Best Practice Guidelines}) \\
LCB     & Centro longitudinal de carena (\textit{Longitudinal Centre of Buoyancy}) \\
KRISO   & \textit{Korea Research Institute of Ships and Ocean Engineering} \\
KCS     & \textit{KRISO Container Ship} \\
KVLCC2  & \textit{KRISO Very Large Crude Carrier}, versión 2 \\
LabHiNO & Laboratorio de Hidrodinámica Naval y Oceánica (UBA) \\
UBA     & Universidad de Buenos Aires \\
MARIN   & \textit{Maritime Research Institute Netherlands} \\
ASME    & \textit{American Society of Mechanical Engineers} \\
IMO     & \textit{International Maritime Organization} \\
\end{longtable}

\renewcommand{\arraystretch}{1.0}